Software systems tend to grow in complexity and size as they evolve, a phenomenon that Lehman formalized in his Laws of Software Evolution~\cite{Lehman1980Laws}: actively used software must continually adapt to its environment, and that adaptation inevitably increased complexity unless deliberate effort was made to reduce it. Niklaus Wirth revisited this concern in 1995, arguing that the trend had persisted and that software's resource demands continued to outpace its functional improvements~\cite{Wirth1995Plea}. This long-term growth poses challenges to the core principles of the GNU Project philosophy, which emphasises the ability of users to read, understand, and extend the software they run~\cite{Sherlock2018OpenSS}. Empirical studies of software maintenance have further reported that increased complexity was associated with measurable increases in maintenance cost and defect density~\cite{nguyen2017impact}. Debian, one of the oldest and most widely used Linux distributions, is the basis for numerous derivative distributions, including Ubuntu; according to W3Techs, Debian and its derivatives collectively accounted for over 40\% of all Linux-powered websites~\cite{W3Techs2024Linux}. Changes in Debian's core packages therefore propagate into a significant share of contemporary computing infrastructure, which motivates a systematic, descriptive account of how those packages have developed since the distribution's earliest stable release.

\section*{Research Gap}
While the general phenomenon of increasing software size has been widely discussed, and longitudinal studies exist for individual software packages~\cite{molnar2020longitudinal}, no systematic descriptive account exists of how Debian's persistent core packages have developed across the distribution's full release history. In particular, no publicly available pipeline characterises installed size, inter-package dependency structure, lines of code, per-function cyclomatic complexity, and distribution-specific patch volume jointly across every stable Debian release. This thesis addresses that gap by building such a pipeline and applying it to every persistent package in the minbase installation across the 13 stable Debian releases from Potato (2000) to Trixie (2025). The work is framed as exploratory research: it describes what can be observed using the chosen set of tools and metrics, without claiming a causal or correlational relationship between the observed dimensions.

\section*{Research Questions}
Following the exploratory-descriptive framing above, this thesis addresses the following primary research question:

\begin{itemize}
    \item What does a reproducible release-wise analysis of the persistent packages in Debian's minbase installation, using an open-source tool pipeline, reveal about their development across all stable releases from Potato (2000) to Trixie (2025), characterised along the dimensions of installed size, dependency structure, lines of code, per-function cyclomatic complexity, and distribution-specific patch volume?
\end{itemize}

Two secondary questions sharpen the primary one:
\begin{itemize}
    \item How has the composition of Debian's minbase installation changed across releases, specifically which packages entered, which were removed, and which form the persistent core?
    \item How has the volume and file-level structure of distribution-specific patches maintained by Debian on top of upstream source code evolved across releases?
\end{itemize}

The research questions are deliberately descriptive. The thesis does not ask whether observed changes in one dimension are correlated with, caused by, or predictive of changes in another; such inferential claims would require larger samples, controlled comparisons, and tools beyond the scope of this work.

\section*{Method}
The methodological mode of this thesis is exploratory-descriptive: a phenomenon (the release-by-release development of Debian's persistent core packages) is characterised using a defined toolchain, and the characterisation itself is the contribution. The method of data collection is a release-wise analysis of every persistent package in the minbase installation, from which five metrics are extracted per release: installed size from the Debian \texttt{Packages} index, dependency edges from the \texttt{Depends} field of that same index, lines of code from the upstream source tarballs using \texttt{tokei}, distribution-specific patch volume from the Debian patch archive, and per-function cyclomatic complexity~\cite{CyclomaticComplexity} from the C and C++ source files of those tarballs using \texttt{pmccabe}. Because cyclomatic complexity is defined at the level of a single function rather than at the level of a package, package- and release-level quantities are always reported as aggregates (mean, median, maximum, sum) over the per-function measurements. The full pipeline is containerised with Docker and orchestrated by a Makefile to make the analysis reproducible. A detailed description of the data collection and analysis methodology is provided in \autoref{ch:method}.

\section*{Expected Output}
The deliverables of this thesis are a reproducible analysis pipeline together with the descriptive artefacts it produces when applied to the 13 stable Debian releases between 2000 and 2025. Concretely, these comprise:

\begin{itemize}
    \item A release-wise timeline of the minbase installation's package composition, reporting the package count of every release, resolving its largest movements package by package, and identifying the 25 source packages that form the persistent core.
    \item Per-release summary statistics and visualisations for installed size (minbase total and persistent subset), dependency structure (per-package dependency counts), lines of code (total and decomposed by programming language), per-function cyclomatic complexity (mean, median, maximum, and per-release distribution), and distribution-specific patch volume.
    \item A containerised pipeline of established open-source tools (\texttt{debootstrap}, \texttt{tokei}, \texttt{pmccabe}, \texttt{gnuplot}) and supporting shell and AWK scripts, released alongside the thesis and intended to be reproducible and extensible by other researchers.
    \item An interactive companion application, the Debian Bloat Explorer, built with the Streamlit framework, that reads the same per-release CSV files as the plotting scripts and supports drill-down into individual packages and releases for verifying and diagnosing the aggregate figures.\footnote{A hosted instance is publicly available at \url{https://debian-bloat-explorer.streamlit.app} (accessed 4 June 2026); source code and data availability are documented in \autoref{app:availability}.}
    \item A catalogue of figures (line charts, boxplots, stacked histograms, bar charts) that make the development of each metric visually accessible without themselves constituting substantive claims about the direction or magnitude of that development. The discussion chapter synthesises these descriptive findings and situates them against prior work; interpretation that goes beyond such synthesis, in particular any causal or correlational claim relating one dimension to another, is deferred to future work.
\end{itemize}

Additional per-package metrics such as the Maintainability Index~\cite{Coleman1994Using} were considered during scoping; they are not part of the current pipeline but represent a natural extension.

\section*{Relevance}
The relevance of this thesis lies in producing, and making reproducible, a descriptive baseline of how Debian's persistent core packages have developed since the distribution's earliest stable release. The pipeline and the resulting per-release artefacts are intended as a foundation: future work that wishes to investigate causal mechanisms behind complexity growth, to evaluate automated software de-bloating techniques~\cite{quach2018debloating}, or to compare Debian with other long-lived distributions can take the metrics and visualisations produced here as its starting point, rather than having to rebuild the data-collection pipeline from scratch.

\section*{Project Structure}
The thesis is structured as follows: the introduction establishes the context, research gap, and the exploratory-descriptive framing that guides the remainder of the work. The related work section reviews existing literature on software complexity metrics, studies of open-source software evolution, and prior analyses of Debian's package management system. The method section describes the data-collection and analysis pipeline in detail, including the tools used, the definition of the minbase installation, the identification of persistent packages, and the exact procedure for computing each metric. The results section presents the per-release artefacts produced by the pipeline: summary statistics, figures, and composition timelines. The discussion section reviews what can and cannot legitimately be said from those artefacts, names the limitations of the approach, and identifies directions for future work. The conclusion summarises the artefacts produced and the scope of the claims they support.